\appendix
\section{Omitted proofs}
\label{omitted-proofs}
\begin{proof}[Proof of \cref{loop-PG}]
  % Note that any element $f:\Aut\P G$, \ie morphism of actions $f:\P G\to\P G$,
  % satisfies $xf(y)=f(xy)$ for any $x,y:G$.
  The two functions
  \begin{align*}
    \phi:\Aut\P G&\to G
    &
    \psi:G&\to\Aut\P G
    \\
    f&\mapsto f(1)
    &
    x&\mapsto y\mapsto yx
  \end{align*}
  % $\phi:\Aut\P G\to G$ and $\psi:G\to\Aut\P G$, respectively defined by
  % $\phi(f)\defd f(1)$ and $\psi(x)(y)\defd yx$
  are group morphisms. Namely, given $f,g:\Aut\P G$, we have
  \begin{align*}
    \phi(g\circ f)&=g\circ f(1)=g(f(1)1)=f(1)g(1)=\phi(f)\phi(g)
    &
    \phi(\id)&=\id(1)=1
  \end{align*}
  and given $x,y:G$, we have for every $z:G$,
  \begin{align*}
    \psi(xy)(z)&=z(xy)=(zx)y=\psi(y)\circ\psi(x)(z)
    &
    \psi(1)(x)&=x1=\id(x)
  \end{align*}
  Moreover, they are mutually inverse. Namely, given $f:\Aut\P G$ and $x:G$, we
  have
  \begin{align*}
    \psi\circ\phi(f)(x)&=xf(1)=f(x1)=f(x)
    &
    \phi\circ\psi(x)&=1x=x    
  \end{align*}
  We thus have $\Aut\P G\simeq G$ and we conclude by univalence.
\end{proof}

\begin{proof}[Proof of \cref{fiberwise-embedding}]
  By definition, the map $\Sigma A.f$ is an embedding iff for every $(a,x)$ and
  $(a',x')$ in $\Sigma A X$, the induced map
  \[
    (a, x) = (a', x') \to (a, f\,a\,x) = (a', f\,a'\,x')
  \]
  is an equivalence. By the characterization of equalities in
  $\Sigma$-types~\cite[Theorem~2.7.2]{hottbook}, this map corresponds to a map
  \[
    (\Sigma(p : a = a').\subst Pp(x) = x') \to (\Sigma(p : a = a').\subst Qp(f\,a\,x) = f\,a'\,x')
  \]
  By \cite[Theorem~4.7.7]{hottbook}, this is an equivalence if and only if the
  fiber map
  \[
    (\subst Pp(x) = x') \to (\subst Qp(f\,a\,x) = f\,a'\,x')
  \]
  is an equivalence for every $p : a = a'$. By path induction, this is true if
  and only if
  \[
    \ap{(f\,a)} : x = x' \to f\,a\,x = f\,a\,x'
  \]
  is an equivalence for all $a : A$, and $x,x' : X a$. By definition, this is
  the requirement that $f\,a$ is an embedding for all $a : A$.
\end{proof}

\section{Equivalence between internal and external groups}
\paragraph{Functoriality of delooping}
\label{delooping-functorial}
One of the main properties of the delooping operation is that it is a ``local
inverse'' to taking loop spaces in the following sense:

\begin{proposition}
  \label{delooping-morphisms}
  Given pointed connected groupoids $A$ and~$B$ and a group morphism
  $f:\Loop A\to\Loop B$, there is a unique pointed morphism $g:A\to B$ such
  that $\Loop g=f$.
\end{proposition}
\begin{proof}
  By~\cite[Corollary~12]{warn2023eilenberg}, with $n=0$, we have that the type
  $\Sigma(g:A\to B).\Loop g=f$ is equivalent to
  $(p,q:\pt{A}=\pt{A})\to f(p\pcomp q)=f(p)\pcomp f(q)$. This type is a
  proposition because $\Loop B$ is a set (because $B$ is a groupoid), and
  inhabited (because $f$ is a morphism of groups), and thus contractible.
\end{proof}

\noindent
As an immediate consequence of the above lemma, deloopings are unique:

\begin{proposition}
  \label{delooping-unique}
  Given two deloopings $\B G$ and $\B' G$ of a group~$G$, we have $\B G=\B' G$.
\end{proposition}

\noindent
Given a group morphism $f:G\to H$ such that both~$G$ and~$H$ admit deloopings
(and this actually always holds by \cref{PG-delooping}), the \emph{delooping} of~$f$ is the morphism
\[
  \B f:\B G\to\B H
\]
associated, by \cref{delooping-morphisms}, to the morphism
$\transportinv{\Beq H}\circ f\circ\transport{\Beq G}:\Loop\B G\to\Loop\B
H$. By \cref{delooping-morphisms}, this operation is functorial in the sense
that that it preserves identities and composition.

\paragraph{Equivalence between the two points of view}
Although this is not central in this article, we shall mention here the
fundamental equivalence provided by the above constructions; details can be
found in~\cite{symmetry}.
%
We write $\IGroup$ for the type of internal groups, \ie pointed connected
groupoids.

% fact that the functions $\Loop$ and $\B$ defined above actually induce an
% equivalence between pointed connected groupoids and groups.

\begin{theorem}
  \label{internal-external-equivalence}
  The maps $\Loop:\IGroup\to\Group$ and $\B:\Group\to\IGroup$ form an
  equivalence of types.
\end{theorem}
\begin{proof}
  Given a group~$G$, we have $\Loop\B G=G$ by definition of~$\B G$. Given an
  internal group~$A$, we have $\B\Loop A\simeq A$ by \cref{delooping-unique}.
\end{proof}

\noindent
The above theorem thus states looping and delooping operators allow us to go
back and forth between the external and the internal point of view of group
theory in homotopy type theory. Note that the torsor construction only gives a
delooping in a larger universe than the original group unless one makes
additional assumptions such as the \emph{replacement
  axiom}~\cite[Axiom~18.1.8]{rijke2022introduction}.

% \noindent
% More generally, it is expected that the types $\IGroup[n]$ of pointed connected
% $n$-types provide a reasonable notion of ``higher-dimensional groups''.

\paragraph{Internal group actions}
In a similar way that the traditional notion of group admits an internal
reformulation (\cref{internal-groups}), the notion of action also admits an
internal counterpart which can be defined as follows.
%
Given a group~$G$, an \emph{internal} action of $G$ on a set $A$ is a function
\[
  \alpha : \B G \to \Set
\]
such that $\alpha(\pt{})=A$. Since $\Set$ is a
groupoid~\cite[Theorem 7.1.11]{hottbook}, by
\cref{internal-external-equivalence}, we have equivalences of types
\[
  (\B G\to\Set)
  \quad\equivto\quad
  (\Loop\B G\to\Loop(\Set,A))
  \quad\equivto\quad
  (G\to\Aut A)
\]
which show that internal group actions correspond to external ones: the
delooping operator internalizes an external group action, and the looping
operator externalizes an internal group action.
% \EO{faire le même schema pour les groupes internes et externes avec $\B$ et
% $\Loop$ d'un coté et de l'autre}


% BEGIN unchanged author comp.tex
\section{Connected components}
\label{connected-component}
We define the \emph{connected component} of the pointed type~$A$ as the type of
points which are merely connected to the distinguished point of~$A$. This type
is noted~$\Comp A$ (or $\Comp(A,\pt{})$ when we want to specify the
distinguished element~$\pt{}$). Formally,
\[
  \Comp A
  \qdefd
  \Sigma(x:A).\ptrunc{\pt{}=x}
\]
This type is canonically pointed by $(\pt{},\ptrunq{\refl{}})$.
%
This construction deserves its name because it produces a connected space, whose
geometry is the same as the original space around the distinguished point, as
shown in the following two lemmas.

\begin{lemma}[\agda{Comp}{isConnectedComp}]
  The type $\Comp A$ is connected.
\end{lemma}
\begin{proof}
  It can be shown that a type~$X$ is connected precisely when both $\ptrunc{X}$
  and $(x,y:X)\to\ptrunc{x=y}$ are inhabited, \ie when~$X$ merely has a point
  and any two points are merely equal~\cite[Exercise~7.6]{hottbook}.
  %
  In our case, the type $\Comp A$ is pointed and thus $\ptrunc{\Comp A}$
  holds.
  %
  Moreover, suppose that there are two points $(x,p)$ and $(y,q)$ in $\Comp A$
  with $x,y:A$, $p:\ptrunc{\pt{}=x}$ and $q:\ptrunc{\pt{}=y}$. Our goal is to
  show that $\ptrunc{(x,p)=(y,q)}$ holds, which is a proposition, so by
  elimination of propositional truncation, we can therefore assume that $p$
  (\resp $q$) has type $\pt{}=x$ (\resp $\pt{}=y$). Hence, we can construct a path
  $\sym p\pcomp q$ of type $x=y$, and therefore $(x,p)=(y,q)$ because the second
  components belong to a proposition by propositional truncation. We conclude
  that $\ptrunc{(x,p)=(y,q)}$ and finally that $\Comp A$ is connected.
\end{proof}

\begin{lemma}[\agda{Comp}{loopCompIsLoop}]
  \label{loop-comp}
  We have $\Loop\Comp A=\Loop A$.
\end{lemma}
\begin{proof}
  % we want to encode *=x in Comp A as B x = (*=x) and therefore it is enough to
  % show that Sigma A.B is contractible
  We begin by showing that the type
  \begin{equation}
    \label{eq:loop-comp-contr}
    \Sigma((x,t):\Comp A).(\pt{}=x)
  \end{equation}
  is contractible. In order to do so, observe that we have the following
  equivalence of types:
  \begin{align*}
    \Sigma((x,t):\Comp A).(\pt{}=x)
    &\simeq
    \Sigma((x,t):\Sigma(x:A).\ptrunc{\pt{}=x}).(\pt{}=x)
    \\
    &\equivto
    \Sigma(x:A).(\ptrunc{\pt{}=x})\times(\pt{}=x)
    \\
    &\equivto
    \Sigma((x,p):\Sigma(x:A).(\pt{}=x)).\ptrunc{\pt{}=x}
  \end{align*}
  using classical associativity and commutativity properties of
  $\Sigma$-types. Moreover, the type $\Sigma(x:A).(\pt{}=x)$ is
  contractible~\cite[Lemma~3.11.8]{hottbook}, therefore the whole type on the
  last line is a proposition (as a sum of propositions over a proposition), and
  therefore also the original type \cref{eq:loop-comp-contr}. We write $\pt{}'$
  for the element $(\pt{},\ptrunc{\refl{\pt{}}})$ of $\Comp A$. The type
  \cref{eq:loop-comp-contr} is pointed by the canonical element
  $(\pt{}',\refl{})$ and thus contractible as a pointed proposition.

  We have a morphism
  \[
    \begin{array}{r@{\ }c@{\ }c@{\ }c@{\ }c@{\ }l}
      F:&((x,t):\Comp A)&\to&(\pt{}'=(x,t))&\to&(\pt{}=x)\\
      &(x,t)&&p&\mapsto&\cong\fst(p)
    \end{array}
  \]
  sending a path~$p$ to the path obtained by applying the first projection. It
  canonically induces a morphism
  \begin{align*}
    \Sigma((x,t):\Comp A).(\pt{}'=(x,t))&\to\Sigma((x,t):\Comp A).(\pt{}=x)\\
    ((x,t),p)&\mapsto((x,t),\cong\fst(p))
  \end{align*}
  between the corresponding total spaces. Since the left member is contractible
  (by~\cite[Lemma~3.11.8]{hottbook} again) and the right member is contractible
  (as shown above), this is an equivalence. By~\cite[Theorem~4.7.7]{hottbook},
  for every $x:\Comp A$, the fiber morphism $F x$ is also an equivalence. In
  particular, with $x$ being $\pt{}'$, we obtain $\Loop\Comp A\simeq\Loop A$ (as
  a type) and we can conclude by univalence. Note that the equivalence preserves
  the group structure so that the equality also holds in groups.
\end{proof}

\noindent
As a direct corollary of the two above lemmas, we have:

\begin{proposition}
  \label{B-loop}
  Given a pointed groupoid~$A$, $\Comp A$ is a delooping of~$\Loop A$.
\end{proposition}

\begin{remark}
  Some people write $\Aut A$ for $\Loop A$ and the above proposition states that
  we have $\B\Aut A=\Comp A$. For this reason, the (confusing) notation
  $\operatorname{BAut}A$ is also found in the literature for $\Comp A$.
\end{remark}

% Local Variables:
% TeX-master: "article"
% End:

% END unchanged author comp.tex


\section{The flattening lemma}
We recall here the classical flattening lemma, see \cite[Section~6.12]{hottbook}
for a more detailed presentation and proof.

\begin{lemma}[Flattening for coequalizers]
  \label{flattening-coequalizers}
  Suppose given a coequalizer
  \[
    \begin{tikzcd}
      A\ar[r,shift left,"f"]\ar[r,shift right,"g"']&B\ar[r,dotted,"h"]&C
    \end{tikzcd}
  \]
  with $p:h\circ f=h\circ g$, and a type family $P:C\to\U$. Then the diagram
  \[
    \begin{tikzcd}[sep=huge]
      \Sigma A.(P\circ h\circ f)\ar[r,shift left,"\Sigma f.(\lambda\_.\id{})"]\ar[r,shift right,"\Sigma g.e"']&\Sigma B.(P\circ h)\ar[r,dotted,"\Sigma h.(\lambda\_.\id{})"]&\Sigma C.P
    \end{tikzcd}
  \]
  is a coequalizer, where the map
  \[
    e:(a:A)\to P(h(f(a)))\to P(h(g(a)))
  \]
  is induced by transport along~$p$ by
  \[
    e\,a\,x\defd
    \subst{P}{\happly pa}(x)
    \text.
  \]
\end{lemma}

\noindent
Note that there is a slight asymmetry: we could have formulated a similar
statement with $\Sigma A.(P\circ h\circ g)$ as left object.

% \section{Rijke's proof}
% \input{rijke}

% Local Variables:
% TeX-master: "article"
% End:
